\section{Characters of Finite Abelian Groups}
Let $G$ be a finite abelian group with the identity element $1_G$. Let $S^1$ be the set of complex numbers with absolute value $1$. Then consider the multiplicative group $S^1$ with the operation of multiplication.
\begin{Definition}{Character $\chi$ of $G$}{}
  \index{character}%
  A \emph{character} $\chi$ of $G$ is a group homomorphism $\chi:G\to S^1$ to the multiplicative group $S^1$ i.e. for all $g_1,g_2\in G$ $$\chi(g_1\cdot g_2)=\chi(g_1)\cdot \chi(g_2)$$
\end{Definition}

So if $\chi$ is any character of $G$ then $\chi(1_G)=1$. Since by definition for all $g\in G$, $$\chi(g)^{|G|}=\chi(g^{|G|})=\chi(1_G)=1$$ the values of $\chi$ are $|G|^{th}$ roots of unity. Also we have for any $g\in G$, $$\chi(g)\chi(g^{-1})=\chi(g\cdot g^{-1})=\chi(1_G)=1$$ so $\chi(g^{-1})=\chi(g)^{-1}=\overline{\chi(g)}$.

There are some special characters of $G$ which have specific name. The character $\chi_0$ defined by $\chi_0(g)=1$ for all $g\in G$ is called the \emph{trivial character}\index{character!trivial}. For any character $\chi$ of $G$, the character $\overline{\chi}$ defined by $\overline{\chi}(g)=\overline{\chi(g)}$ for all $g\in G$ is called the \emph{conjugate character}\index{character!conjugate} of $\chi$.

Given any finitely many characters of $G$, $\chi_1,\dots, \chi_n$ we can define the product character $\chi=\chi_1\cdot \chi_2\cdots \chi_n$ such that for all $g\in G$, $$\chi(g)=\chi_1(g)\cdot \chi_2(g)\cdots \chi_n(g)$$
Hence for all $g\in G$, $$\chi(g)\cdot \overline{\chi}(g)=|\chi(g)|^2=1$$Hence the $\overline{\chi}$ is also the inverse character of $\chi$. Hence the set of all characters of $G$ forms a group, $\Hom(G,S^1)$ under multiplication whose identity element is the trivial character $\chi_0$.
\begin{Definition}{Dual Group to $G$}{}
  \index{dual group}%
  For a finite abelian group $G$, the group of all characters of $G$ with the identity element $\chi_0$ is called the \emph{dual group} to $G$ and is denoted by $\widehat{G}$ i.e. $\hG\coloneqq \Hom(G,S^1)$.
\end{Definition}

\begin{observation}
  The dual group $\widehat{G}$ is a finite abelian group. So we can talk about characters of $\hG$ too.
\end{observation}
Since the values of characters of $G$ can only be $|G|^{th}$ roots of unity, the dual group $\widehat{G}$ is a finite abelian group.
\begin{lemma}{Charaters of Cyclic Group}{lem:chars-cyclic-group}
  If $G$ is a finite cyclic group of order $n$ and let $g$ be the generator of $G$, then for all $j\in\{0,1,\dots, n-1\}$ define the function $$\chi_j(g^k)=e^{2\pi i\frac{j\cdot k}{n}} \quad \forall\ k\in \{0,1,\dots, n-1\}$$Then $\hG=\{\chi_j\colon j\in\{0,1,\dots, n-1\}$.
\end{lemma}
\begin{proof}
  Certainly for all $j\in\{0,1,\dots, n-1\}$ defines a character of $G$. On the other hand suppose, $\chi\in\hG$. Then $\chi(g)$ must be an $n^{th}$ root of unity. Hence there exists $j\in\{0,1,\dots, n-1\}$ such that $\chi(g)=e(\nicefrac{j}{n})$. Therefore $\chi=\chi_j$. Hence $\hG=\{\chi_j\colon  0\leq j\leq n-1\}$.
\end{proof}
\begin{observation}
  The characters defined in above lemma forms a cyclic group with $\chi_1$ being the generator of $\hG$.
\end{observation}
\begin{note}
  When $G=(\bbZ/n\bbZ)^*$, the characters of $G$ described by $\chi_j(\alpha)=e^{2\pi i\frac{\alpha\cdot j}{n}}$ like in \lemref{lem:chars-cyclic-group} are precisely the classical \emph{Dirichlet characters}\index{Dirichlet character} modulo $n$, which are of fundamental importance in analytic number theory.
\end{note}
\begin{corolary}{Isomorphism of a Cyclic Group with its Dual}{}
  If $G$ is a finite cyclic group of order $n$ with $g$ being the generator then $G\cong \hG$ by the isomorphism $\varphi:g^j\mapsto \chi_j$ for all $j\in\{0,1,\dots, n-1\}$.
\end{corolary}

Now we will show that we can extend a character of any subgroup to a character of the group.
\begin{Theorem}{Extending Characters from Subgroup to Group}{thm:subgroup-char-extend}
  Let $H$ be a subgroup of the finite abelian group $G$ and let $\chi_H$ be a character of $H$. Then $\chi_H$ can be extended to a character of $G$ that is there exists a character $\chi\in\hG$ with $\chi(h)=\chi_H(h)$ for all $h\in H$.
\end{Theorem}
\begin{proof}
  Let $H$ be a proper subgroup of $G$. Let $\alpha\in G\setminus H$. Then let $H_1$ be the group generated by $H$ and $\alpha$ i.e. $H_1=\langle H,a\rangle$. Let $m$ be the smallest positive integer such that $\alpha^m\in H$. Then for all $g\in H_1$, $g$ can be written uniquely as $g=\alpha^j\cdot h$ where $0\leq j<m$ and $h\in H$.

  Let $\om$ be a complex number such that $\om^m=\chi_H(\alpha^m)$. Then for all $g\in H_1$ where $g=\alpha^j\cdot h$, $h\in H$ define the function $\chi_1(g)=\om^j \chi_H(h)$. Then for any $g_1,g_2\in H_1$ where $g_1=\alpha^i\cdot h_1$, $g_2=\alpha^j\cdot h_2$, we have $g_1\cdot g_2=\alpha^{j+k}\cdot h_1\cdot h_2=\alpha^{j+k\bmod{m}}\cdot h_1\cdot h_2$. Hence $$\chi_1(g_1\cdot \chi_2)=\om^{j+k \bmod{m}}\chi_H(h_1)\cdot \chi_H(h_2)=\alpha^j\chi_H(h_1)\cdot \om^k\chi_H(h_2)=\chi_1(g_1)\cdot \chi_1(g_2)$$Hence $\chi_1$ is a character of $H_1$. And also for all $h\in H$ we have $\chi_1(h)=\chi_H(h)$.

  So if $H_1=G$ we are done. Otherwise we can continue this process above until we get $G$. This process we have to do finitely many times as $G$ is a finite group.
\end{proof}


Therefore if we have two distinct elements $g_1,g_2\in G$, $g_1\neq g_2$ then take the subgroup generated by the element $h=g_1\cdot g_2^{-1}\neq 1$. Since this subgroup is a cyclic subgroup generated by $h$, by \lemref{lem:chars-cyclic-group} we get a character $\chi'$ of the cyclic subgroup such that $\chi'(h)\neq 1$. Then by \thmref{thm:subgroup-char-extend} we get a character $\chi$ of $G$ such that $\chi(h)=\chi'(h)\neq 1$ and hence $\chi(g_1)\neq \chi(g_2)$. Therefore we have the following observation
\begin{observation}
  For any two distinct elements $g_1,g_2\in G$, $g_1\neq g_2$ we can get a character $\chi\in\hG$ such that $\chi(g_1)\neq \chi(g_2)$.
\end{observation}

Now consider any character $\chi$ of $G$. If $g\in G$ such that $\chi(g)=1$ then $\chi(g^{-1})=1$. Hence the set of all elements $g\in G$ such that $\chi(g)$ forms a group. This group is called the \emph{annihilator} of $\chi$ and is denoted by $\Ann(\chi)$.

Similarly suppose $H$ is a subgroup of $G$. Then if $\chi\in\hG$ is a character of $G$ such that $\chi(h)=1$ for all $h\in H$ then $\ov{\chi}(h)=1$ for all $h\in H$. Hence the set of all characters $\chi\in \hG$ such that $\chi(h)=1$ for all $h\in H$ forms a group. This group is called the \emph{annihilator} of $H$ and is denoted by $\Ann(H)$. Since in both cases the notation $\Ann$ is used usually the context will make it clear whether we are talking about the annihilator of a character or the annihilator of a subgroup.

\begin{Definition}{Annihilator}{}
  \index{annihilator!of a character}\index{annihilator!of a subgroup}%
  Let $G$ be a finite abelian group. For any character $\chi\in \hG$, the \emph{annihilator} of $\chi$, $\Ann(\chi)=\{g\in G\colon \chi(g)=1\}$ which is a subgroup of $G$. \setlength{\parindent}{15pt}

  For any subgroup $H$ of $G$, the \emph{annihilator} of $H$, $\Ann(H)=\{\chi\in \hG\colon \chi(h)=1 \text{ for all } h\in H\}$ which is a subgroup of $\hG$.
\end{Definition}


Now every element of $\hG$ is a function $\chi$ which takes an element of $G$ and maps to a complex number in $S^1$. Similarly we can think of every element of as a function from $\hG$ to $S^1$. In fact for every element $g\in G$ define the homomorphism $\ev_g:\hG\to S^1$ where $\ev_g(\chi)=\chi(g)$. Hence $\ev_g$ is a character of $\hG$ for all $g\in G$.
\begin{Theorem}{}{}
  Let $G$ be a finite abelian group.   \begin{enumerate}[label=(\roman*)]
    \item If $\chi\in \hG$ be any nontrivial character then \begin{equation}
            \sum_{g\in G}\chi(g)=0
          \end{equation}
    \item If $g\in G$ with $g\neq 1_G$ then \begin{equation}
            \sum_{\chi\in \hG}\chi(g)=0
          \end{equation}
  \end{enumerate}
\end{Theorem}
\begin{proof}
  Since $\chi$ is nontrivial there exists $h\in G$ such that $\chi(h)\neq 1$. Then $$\chi(h)\sum_{g\in G}\chi(g)=\sum_{g\in G}\chi(h\cdot g)= \sum_{g\in G}\chi(g)$$Hence we have $$(\chi(h)-1)\sum_{g\in G}\chi(g)=0$$Since we know $\chi(h)\neq 1$ we have $\sum_{g\in G}\chi(g)=0$.

  Now from the discussion above for any $g\in G$ and $\chi\in hG$ we have $\chi(g)=\ev_g(\chi)$. Hence $$\sum_{\chi\in \hG}\chi(g)=\sum_{\chi\in \hG}\ev_g(\chi)$$Since $g\neq 1_G$, there is a $\chi\in\hG$ such that $\chi(g)\neq 1$. Hence $\ev_g$ is a nontrivial character of $\hG$. Hence by the first part of the theorem we have $\sum_{\chi\in \hG}\ev_g(\chi)=0$.
\end{proof}

The above theorem leaves the case when $\chi$ is the trivial character. Then for all $g\in G$, $\chi(g)=1$. Hence we have $\sum_{g\in G}\chi(g)=|G|$. Hence we have the following corollary
\begin{observation}\label{obs:char-sum}
  For any finite abelian group $G$. Then \begin{enumerate}[label=(\roman*)]
    \item Let $\chi$ is a character of $G$. \begin{equation}\label{eq:char-sum}
            \sum_{g\in G}\chi(g)=\begin{cases}
              |G| & \text{if $\chi=\chi_0$} \\
              0   & \text{otherwise}
            \end{cases}
          \end{equation}
    \item $|G|=\sum_{\chi\in \hG}\sum_{g\in G}\chi(g)$
  \end{enumerate}
\end{observation}

So now by counting in two ways for the sum $\sum_{\chi\in \hG}\sum_{g\in G}\chi(g)$ we get the following theorem:
\begin{Theorem}{}{thm:group-char-equal}
  For any finite abelian group $G$, we have $|G|=|\hG|$.
\end{Theorem}
\begin{proof}
  From \obsref[(ii)]{obs:char-sum} we have $$|G|=\sum_{\chi\in \hG}\sum_{g\in G}\chi(g)=\sum_{g\in G}\sum_{\chi\in \hG}\chi(g)=|\hG|$$Hence we have the result.
\end{proof}

From above theorem and \obsref[(i)]{obs:char-sum} we get the \emph{orthogonality relations for characters}.\index{orthogonality relations!for characters}
\begin{observation}\label{obs:orthogonal-relation}
  Let $G$ be a finite abelian group.
  \begin{enumerate}[label=(\roman*)]
    \item Let $\chi,\tau$ are characters of $G$. Then $$\frac{1}{|G|}\sum_{g\in G}\chi(g)\cdot \overline {\tau(g)}=\begin{cases}
              1 & \text{if $\chi=\tau$} \\
              0 & \text{otherwise}
            \end{cases}$$
    \item If $g,h\in G$ are any two elements then $$\frac{1}{|G|}\sum_{\chi\in \hG}\chi(g)\cdot \overline{\chi(h)}=\begin{cases}
              1 & \text{if $g=h$}  \\
              0 & \text{otherwise}
            \end{cases}$$
  \end{enumerate}
\end{observation}

Since we already showed that for every element $g\in G$, $\ev_g$ is a character of $\hG$. \thmref{thm:group-char-equal} gives us the possibility that $G$ and $\hG$ are isomorphic. In \lemref{lem:chars-cyclic-group} we already showed that for cyclic groups $G$ and $\hG$ are isomorphic. In the next few results we will show that this is the case for all finite abelian groups.

By Fundamental Theory of Finite Abelian Groups every finite abelian group $G$ can be written as a direct product of finitely many cyclic groups. Hence it is enough to show that if $G_1$ and $G_2$ are two finite abelian groups and $G=G_1\times G_2$ then $\hG$ is isomorphic to the direct product of $\hG_1$ and $\hG_2$.
\begin{Theorem}{Isomorphism of Product of Dual Groups with Dual of Their Product}{thm:group-product-char}
  Let $G_1$ and $G_2$ be two finite abelian groups and let $G=G_1\times G_2$. Then $\hG\cong \hG_1\times \hG_2$.
\end{Theorem}
\begin{proof}
  Since $G$ is the direct product of $G_1$ and $G_2$, $G$ consists of pairs $(g_1,g_2)$ where $g_1\in G_1$ and $g_2\in G_2$. So consider the map $\varphi:\hG_1\times \hG_2\to \hG$ where for any $\chi_1\in \hG_1$ and $\chi_2\in \hG_2$, $\chi\coloneqq \varphi(\chi_1,\chi_2)$ is the character of $G$ defined by the following: $$\text{for all }(g_1,g_2)\in G\qquad \chi(g_1,g_2)=\chi(g_1,1)\cdot \chi(1,g_2)=\chi_1(g_1)\cdot \chi_2(g_2)$$Hence $\varphi$ is a homomorphism. Now we will show that $\varphi$ is an isomorphism.

  Let $\chi\in \hG$. Then define $\chi_1(g_1)\coloneqq\chi(g_1,1)$ for all $g_1\in G_1$ and similarly define $\chi_2(g_2)\coloneqq\chi(1,g_2)$ for all $g_2\in G_2$. Then $\chi_1$ is a character of $G_1$ and $\chi_2$ is a character of $G_2$. And also for any $(g_1,g_2)\in G$ we have $$\chi_1(g_1)\cdot \chi_2(g_2)=\chi(g_1,1)\cdot \chi(1,g_2)=\chi(g_1,g_2)$$ Hence $\varphi(\chi_1,\chi_2)=\chi$. Hence $\varphi$ is an isomorphism. Therefore $\hG\cong \hG_1\times \hG_2$.
\end{proof}

\begin{Theorem}{Isomorphism of Group with its Dual}{thm:group-char-isomorphism}
  For any finite abelian group $G$, $G$ is isomorphic to $\hG$.
\end{Theorem}
\begin{proof}
  By Fundamental Theorem of Finite Abelian Groups, $G$ can be written as a direct product of finitely many cyclic groups i.e. there exists cyclic groups $Z_{n_1},Z_{n_2},\dots, Z_{n_k}$ such that $$G\cong Z_{n_1}\times Z_{n_2}\times \cdots \times Z_{n_k}$$ Hence by repeated use of \thmref{thm:group-product-char} we have $\hG\cong \hZ_{n_1}\times \hZ_{n_2}\times \cdots \times \hZ_{n_k}$. Since by \lemref{lem:chars-cyclic-group} for all $i\in\{1,2,\dots, k\}$, $\hZ_{n_i}\cong Z_{n_i}$ we have $\hG\cong Z_{n_1}\times Z_{n_2}\times \cdots \times Z_{n_k}\cong G$ Hence we have the result.
\end{proof}